\chapter*{\textcolor{gray}{Introduction}}
Le cerveau, merveille de la nature, réalise des computations complexes pour transformer les données sensorielles en perceptions et actions. Ces processus, opérant à l'échelle de la milliseconde, reposent sur des mécanismes biophysiques au niveau des neurones et de leurs composants. Loin des modèles simplistes des réseaux de neurones, ces derniers effectuent des opérations sophistiquées en exploitant des variables comme le potentiel de membrane ou les concentrations en ions sodium et potassium. Mais comment fonctionnent ces "circuits biologiques" ? Quels sont leurs éléments fondamentaux ? Cette présentation explore les principes biophysiques sous-jacents à ces computations, du neurone complet aux modèles de la membrane neuronale (chapitre 2 et 3), en passant par le modèle emblématique de Hodgkin-Huxley (chapitre 4) et les analyses mathématiques et numériques (chapitres 5 et 6). Plongeons dans l'univers fascinant de la computation neuronale !